\subsection{The exact reduced nonproperness locus}
\label{subsec:nonproperness}

We now make the fullness condition intrinsic.  Work geometrically over
\(\overline K\).  For a polynomial map \(F:\A^n_{\overline K}\to
\A^n_{\overline K}\), let \(S_F\) be the set of target points at which \(F\)
is not finite over any Zariski neighborhood.  Since source and target are
affine, this is the same as failure of properness over every neighborhood.
Equivalently, compactify the graph in
\(\mathbb P^n_{\overline K}\times\A^n_{\overline K}\); then \(S_F\) is the
image of its boundary.  We give \(S_F\) its reduced closed-subscheme
structure.  Over \(\mathbb C\) this agrees with the usual analytic Jelonek
set, defined by sequences escaping to infinity.

For a polynomial \(h\), write \(V_{\mathrm{red}}(h)\) for the reduced closed
subscheme \(V(\sqrt{(h)})\).  Put
\begin{equation}
 \Delta_G(\Pi,B,C)
 =\operatorname{Disc}_S\bigl(E_{\Pi,B,C}(S)\bigr).
 \label{eq:gauge-discriminant}
\end{equation}
Here the discriminant is the universal degree-\(N\) polynomial expression in
the coefficients of \(E_{\Pi,B,C}\).  We do not replace it on a special
slice by the discriminant of the lower-degree specialization.

\begin{lemma}[Fiber cardinality and nonproperness in characteristic zero]
\label{lem:properness-cardinality}
Let \(k\) be an algebraically closed field of characteristic zero and let
\(F:\A^n_k\to\A^n_k\) be a dominant \'etale polynomial map of geometric
degree \(d\).  Then
\begin{equation}
 y\notin S_F\quad\Longleftrightarrow\quad
 \#F^{-1}(y)=d.
 \label{eq:properness-cardinality}
\end{equation}
\end{lemma}

\begin{proof}
Jelonek proves this multiplicity criterion for complex polynomial maps
\cite[Proposition~6]{Jelonek1993}.  We record why that stated complex result
is sufficient here.  Descend the finitely many coefficients of \(F\) and the
coordinates of \(y\) to a finitely generated field \(k_0/\mathbb Q\), and
choose an embedding \(k_0\hookrightarrow\mathbb C\).  The graph-boundary
description above is compatible with extension of algebraically closed
fields: a boundary fiber is empty before base change exactly when it is empty
after base change.  Thus membership in \(S_F\) is unchanged.  Because \(F\)
is \'etale, its fiber at \(y\) is a finite \'etale \(k_0\)-algebra; its
dimension, hence its number of geometric points after either algebraic
closure, is also unchanged.  Finally, geometric degree is invariant under
this scalar extension.  Proposition~6 over \(\mathbb C\), where every local
multiplicity is one, therefore gives \eqref{eq:properness-cardinality} over
\(k\).
\end{proof}

\begin{theorem}[Exact reduced quadratic-gauge nonproperness]
\label{thm:exact-nonproperness}
For every seed \eqref{eq:seed},
\begin{equation}
 \boxed{S_{F_G}=V_{\mathrm{red}}(\Delta_G).}
 \label{eq:jelonek-discriminant}
\end{equation}
Equivalently, their sets of geometric points satisfy
\(\lvert S_{F_G}\rvert=V_{\overline K}(\Delta_G)\).
More precisely, the following statements hold.
\begin{enumerate}
\item If \(\pi\neq0\) and the multiplicity partition of
      \(E_{\pi,b,c}\) is \((m_1,\ldots,m_s)\), then
      \begin{equation}
       \#F_G^{-1}(\pi,b,c)=\#\{i:m_i=1\},
       \qquad
       N-\#F_G^{-1}(\pi,b,c)=\sum_{m_i\geq2}m_i.
       \label{eq:partition-sheet-loss}
      \end{equation}
\item On \(\Pi=0\), the complete geometric fiber table is
      \begin{equation}
      \begin{array}{c|c|c}
       \text{target stratum}&\#F_G^{-1}(0,B,C)&\text{lost sheets}\\ \hline
       B\neq0,\ BC\neq1&3&N-3\\
       B\neq0,\ BC=1&1&N-1\\
       B=0&1&N-1.
      \end{array}
      \label{eq:pi-zero-fiber-table}
      \end{equation}
\item For \(N=3\), the raw resultant has an extraneous leading-coefficient
      factor:
      \[
       \operatorname{Res}_S(E,E')=\text{\rm unit}\cdot\Pi\Delta_G,
       \qquad
       \Delta_G(0,B,C)=\frac{g_1^4}{4}B^2(1-BC).
      \]
      Hence the entire plane \(\Pi=0\) is not a nonproperness component in
      degree three.
\item For \(N\geq4\), put
      \[
       e_N=N^2-3N-2.
      \]
      Then
      \begin{equation}
       \Delta_G=\Pi^{e_N}D_G,
       \qquad
       D_G(0,B,C)=\kappa_G B^2(1-BC),
       \qquad \kappa_G\in K^*.
       \label{eq:saturated-discriminant}
      \end{equation}
      In particular,
      \[
       S_{F_G}=V_{\mathrm{red}}(\Pi D_G),
      \]
      whose underlying set is \(V(\Pi)\cup V(D_G)\), and the resultant and
      discriminant define the same reduced closed subscheme.
\end{enumerate}
\end{theorem}

\begin{proof}
Because \(F_G\) is \'etale, every geometric fiber is finite and reduced.
Its generic cardinality is \(N\), so
Lemma~\ref{lem:properness-cardinality} applies with \(d=N\).

Suppose first that \(\pi\neq0\).  Theorem~\ref{thm:localized-gauge}
identifies the fiber with
\[
 \operatorname{Spec}
 \bigl((\overline K[S]/(E_{\pi,b,c}))[1/\overline{E_{\pi,b,c}'}]\bigr).
\]
Thus a distinct root contributes one affine point exactly when it is simple.
This proves \eqref{eq:partition-sheet-loss}; it also shows that the fiber has
cardinality \(N\) exactly when \(\Delta_G(\pi,b,c)\neq0\).

It remains to analyze \(\Pi=0\), where the inverse chart used above is not
available.  In the source, \(\Pi=tq\).  On \(q=0\), one has
\[
 B=y,\qquad C=2S-BS^2,\qquad S=\frac{x}{t}.
\]
The residual inverse equation is therefore
\begin{equation}
 BS^2-2S+C=0.
 \label{eq:q-zero-residual}
\end{equation}
A root reconstructs on \(q=0\) precisely when \(1-BS\neq0\).  Hence this
divisor contributes two points when \(B\neq0\) and \(BC\neq1\), no point when
\(B\neq0\) and \(BC=1\), and one point when \(B=0\).

On \(t=0\), the relation \(1+xy=0\) gives \(y=-1/x\), while
\[
 q=\frac{g_1}{g_3}y^2,\qquad B=-2y=\frac2x.
\]
Thus \(t=0\) contributes one point when \(B\neq0\) and none when \(B=0\).
Once \(x\) and \(y\) are fixed, the equation for \(C\) is linear in \(z\)
with nonzero coefficient, so the point is unique.  The divisors \(t=0\) and
\(q=0\) do not meet.  This proves \eqref{eq:pi-zero-fiber-table}.

For \(N=3\), specializing the degree-three discriminant at \(\Pi=0\) gives
\[
 \Delta_G(0,B,C)=\frac{g_1^4}{4}B^2(1-BC).
\]
Its zero set is exactly the two short-fiber strata in
\eqref{eq:pi-zero-fiber-table}.  Since the leading coefficient is
\(g_3\Pi\), the resultant has one additional factor \(\Pi\).

Now let \(N\geq4\), put \(r=N-3\), and set
\(R=\overline K[B,C]\).  Over
\(\operatorname{Frac}(R)[[\Pi]]\), the lower Newton polygon of \(E\) has
vertices
\[
 (0,0),\quad(2,0),\quad(3,1),\quad(N,N).
\]
Consequently its root valuations, with multiplicities, are
\[
 0\ (2),\qquad -1\ (1),\qquad
 -\frac{r+2}{r}\ (r).
\]
The first two blocks are exactly the affine \(q=0\) and \(t=0\) branches
found above; the final block consists of the \(N-3\) missing sheets.
The residual equations of the three blocks are, respectively,
\[
 BZ^2-2Z+C=0,\qquad
 Z=\frac{g_1B}{2g_3},\qquad
 g_3+g_NZ^r=0.
\]
They are separable over \(\operatorname{Frac}(R)\): the first has
discriminant \(4(1-BC)\), the second is linear, and the last is separable
because the characteristic is zero and \(g_3g_N\neq0\).

Using
\[
 \operatorname{Disc}(E)
 =\operatorname{lc}(E)^{2N-2}
   \prod_{i<j}(s_i-s_j)^2
\]
and the three displayed root valuations gives
\[
 \operatorname{ord}_{\Pi}(\Delta_G)
 =2N(N-1)-4-4(r+2)-2(r+2)-(r-1)(r+2)
 =N^2-3N-2.
\]
Separability makes this order exact.  To retain the dependence on \(B,C\),
let \(\alpha_1,\alpha_2\) be the roots of the first residual equation, put
\(b_0=g_1B/(2g_3)\), and let \(v_1,\ldots,v_r\) be the roots of
\(g_3+g_NZ^r\).  The leading coefficient in the root-product formula is
\begin{align*}
 &\mathrel{\phantom{=}}\pm g_N^{2N-2}
 (\alpha_1-\alpha_2)^2 b_0^4
 \left(\prod_jv_j\right)^6
 \prod_{i<j}(v_i-v_j)^2\\
 &=\pm\frac{r^r}{4}
 g_1^4g_3^{r+1}g_N^{r-1}B^2(1-BC).
\end{align*}
Indeed,
\[
 (\alpha_1-\alpha_2)^2=\frac{4(1-BC)}{B^2},\quad
 \left(\prod_jv_j\right)^6=\left(\frac{g_3}{g_N}\right)^6,
\quad
 \prod_{i<j}(v_i-v_j)^2
 =\pm r^r\left(\frac{g_3}{g_N}\right)^{r-1}.
\]
These calculations were made in \(\operatorname{Frac}(R)\), but
\(\Pi\)-adic order in the domain \(R[\Pi]\) is unchanged on passing to its
fraction field.  Hence they prove the global polynomial identity
\[
 \Delta_G=\Pi^{e_N}D_G,\qquad
 D_G(0,B,C)=\kappa_GB^2(1-BC),
\quad
 \kappa_G=\pm\frac{r^r}{4}
 g_1^4g_3^{r+1}g_N^{r-1}\neq0.
\]
Thus \(e_N\) is the exact global \(\Pi\)-adic order.  After specializing
\((B,C)\), the pointwise order is \(e_N\) when
\(B(1-BC)\neq0\) and may jump on the two displayed special strata.

No general assertion about discriminants under a drop of leading coefficient
is being used.  For \(N=3\), the explicit specialized formula cuts out
exactly the two short-fiber strata.  For \(N\geq4\), \(\Delta_G\) vanishes on
the whole slice \(\Pi=0\), and every fiber on that slice has fewer than \(N\)
points by \eqref{eq:pi-zero-fiber-table}.  On \(\Pi\neq0\), the ordinary
degree-\(N\) discriminant detects exactly the repeated-root, hence
short-fiber, locus.
The fiber table and \eqref{eq:properness-cardinality} now show that the zero
set of \(\Delta_G\) is exactly the underlying nonproperness set, and taking
radicals gives \eqref{eq:jelonek-discriminant}.
\end{proof}

\begin{proposition}[Boundary normalization and explicit escape arcs]
\label{prop:boundary-arcs}
On \(\Pi\neq0\), the reduced discriminant divisor is birationally
parametrized by
\begin{equation}
 B=\frac{G_\Pi'(s)}{g_1s},
 \qquad
 C=\frac{2G_\Pi(s)-sG_\Pi'(s)}{g_1},
 \qquad (\Pi,s)\in\mathbb G_m^2.
 \label{eq:global-discriminant-parametrization}
\end{equation}
If \(s_0\) is a root of multiplicity \(m\geq2\) over
\((\pi,b,c)\), every one of the \(m\) lost sheets is represented, after a
finite base change, by a Laurent-series source arc.  With root displacement
as uniformizer, its generic valuations are
\[
 \operatorname{ord}(x)=-(m-1),\qquad
 \operatorname{ord}(q)=m-1,\qquad
 \operatorname{ord}(z)\geq m-1.
\]

If \(N\geq4\), put \(d=N-3\) and \(h=\gcd(d,2)\).  At the generic point of
\(\Pi=0\), the \(d\) lost sheets form \(h\) boundary primes, each with
ramification index \(d/h\).  After the common base change \(\Pi=u^d\), their
valuations are
\begin{equation}
\begin{array}{c|rrrrrrr}
 &S&D&t&x&y&q&z\\ \hline
 \operatorname{ord}_u&
 -(N-1)&-(N+1)&N+1&2&-2&-4&-(2N+6).
\end{array}
\label{eq:infinity-valuations}
\end{equation}
For \(N=3\) this generic boundary block is absent.
On the special stratum \(B=0\), two further missing sheets have
\begin{equation}
 \Pi=u^2,\qquad S\sim a u^{-1},\qquad
 a^2=-\frac{g_1}{g_3},
 \qquad
 \bigl(\operatorname{ord}_u x,\operatorname{ord}_u y,
       \operatorname{ord}_u q,\operatorname{ord}_u z\bigr)
 =(-1,1,2,2).
 \label{eq:B-zero-valuations}
\end{equation}
Together with the multiplicity-block arcs above, these account for every
lost sheet on every stratum of \eqref{eq:pi-zero-fiber-table}.
\end{proposition}

\begin{proof}
Equations \eqref{eq:global-discriminant-parametrization} are \(E=E'=0\)
solved for \(B,C\).  Their birationality was proved in the monodromy argument
around \eqref{eq:quadratic-discriminant-parametrization}.

For a repeated root, set \(s(u)=s_0+u\) and keep \(\pi,b\) fixed while putting
\[
 c(u)=\frac{2G_\pi(s(u))}{g_1}-b\,s(u)^2,
 \qquad
 d(u)=\frac{G_\pi'(s(u))}{g_1}-b\,s(u).
\]
Then \(E_{\pi,b,c(u)}(s(u))=0\) identically and
\(\operatorname{ord}_u d=m-1\).  The exact reconstruction formulas
\[
 x=\frac{s}{d},\qquad
 y=b-\beta(\pi,s)-\pi s,\qquad
 q=\pi d,\qquad
 z=\pi d^3-\frac{g_1}{g_3}y^2(d^2+3d)
\]
give the stated valuations.  Applying Newton--Puiseux to the other roots in
the same multiplicity block gives all \(m\) arcs after a finite base change.

Suppose \(N\geq4\).  For \(\Pi=0\), the final Newton edge has horizontal
length \(d\) and endpoint equation
\[
 g_3+g_NZ^d=0.
\]
Equivalently its binomial model is
\[
 W^d=\text{unit}\cdot\Pi^{d+2}.
\]
It has \(h=\gcd(d,2)\) primes, each of ramification index \(d/h\).
Normalizing by
\[
 v(\Pi)=\frac d h,\qquad
 v(S)=-\frac{d+2}{h}
\]
gives
\[
 v(D)=-\frac{d+4}{h},\quad
 v(x)=\frac2h,\quad
 v(y)=-\frac2h,\quad
 v(q)=-\frac4h,\quad
 v(z)=-\frac{2d+12}{h}.
\]
Multiplication by \(h\), corresponding to the common base change
\(\Pi=u^d\), yields \eqref{eq:infinity-valuations}.

If \(B=0\), the first Newton edge changes from
\((0,0)\)--\((2,0)\) to \((1,0)\)--\((3,1)\).  Its two nonzero roots satisfy
\[
 \Pi=u^2,\qquad S\sim a u^{-1},\qquad g_1+g_3a^2=0.
\]
Here \(D\) and \(t=D^{-1}\) have order zero.  Substitution in the exact
reconstruction formulas gives \eqref{eq:B-zero-valuations}.  When \(BC=1\)
with \(B\neq0\), the two additional missing sheets instead form the finite
double-root block already covered by the first part of the proposition.
\end{proof}

\begin{corollary}[Intrinsic fullness of the distinguished fiber]
\label{cor:intrinsic-fullness}
For the seed \(G(S)=P(a+S)-P(a)\) used in
Theorem~\ref{thm:main}, the distinguished target \(y_{P,a}\) lies in the
maximal open subset on which \(\widetilde F_G\) is finite \'etale of degree
\(N\).
\end{corollary}

\begin{proof}
At \(y_{P,a}\), the inverse polynomial is \(P(a+S)\).  Hence
\[
 \Delta_G(y_{P,a})=\operatorname{Disc}(P(a+S))
 =\operatorname{Disc}(P)\neq0.
\]
Theorem~\ref{thm:exact-nonproperness} places the target outside the
nonproperness hypersurface.  Proper and \'etale is finite \'etale, and the
degree is \(N\).
\end{proof}

\begin{remark}[Polynomial curves]
The parametrization \eqref{eq:global-discriminant-parametrization}, with
\(s\) fixed, gives polynomial curves in \(\Pi\) covering the discriminant
component, while the component \(\Pi=0\) is covered by affine lines.  Thus the
explicit geometry realizes directly the general theorem that a complex
Jelonek set is a uniruled hypersurface covered by polynomial curves of bounded
degree \cite{Jelonek1993,JelonekLason}.
\end{remark}
